\subsubsection{The power parameter and the exact fibre convention}

Let $A$ be an integral Banach ring.  For every $x\in\cM(A)$, set
\begin{equation}
 \rho(x):=\frac{\log\max\{1,\abs{2}_x\}}{\log 2}.
\label{eq:rho-base}
\end{equation}
This is a continuous function on $\cM(A)$.

\begin{lemma}
\label{lem:rho-split}
One has $\rho(x)=0$ exactly when the absolute value on the completed residue
field $\mathcal H(x)$ is non-Archimedean.  If $\rho(x)>0$, then
$0<\rho(x)\leq1$, and there is a character into $\C$, unique up to complex
conjugation, such that after completing the fraction field
\begin{equation}
 \abs{a}_x=\abs{\chi_x(a)}_\infty^{\rho(x)}.
\label{eq:character-power}
\end{equation}
Thus $\mathcal H(x)$ is $\R$ or $\C$, with the
$\rho(x)$-th power of the usual absolute value.
\end{lemma}

\begin{proof}
If $\abs{2}_x\leq1$, binary expansion first gives
$\abs{n}_x\leq1$ for all positive integers $n$; applying the binomial
formula to $(a+b)^m$ and taking $m$-th roots gives the ultrametric
inequality.  The converse follows from
$\abs{2}_x\leq\max\{\abs1_x,\abs1_x\}=1$.

If $\abs2_x>1$, the base-expansion proof of Ostrowski's theorem gives
$\abs n_x=n^{\rho(x)}$.  The ordinary triangle inequality gives
$\abs2_x\leq2$, hence $0<\rho(x)\leq1$.  The unpowered function
$v(a)=\abs a_x^{1/\rho(x)}$ is multiplicative, is usual on $\mathbf Q$,
and satisfies the triangle inequality: expanding $(a+b)^m$ yields
\[
 \abs{a+b}_x^m\leq(m+1)
 \bigl(v(a)+v(b)\bigr)^{m\rho(x)},
\]
and the power $1/(m\rho(x))$, followed by $m\to\infty$, proves the
claim.  The completion is a commutative real Banach division algebra, hence
$\R$ or $\C$ by Gelfand--Mazur.  The only isometric ambiguity in the
complex case is conjugation.  This is the character construction used in
\cref{thm:hybrid-weighted-philippon-gauss}.
\end{proof}

Fix henceforth an Archimedean point and abbreviate
$K_\rho=(K,\abs{\cdot}_\infty^\rho)$, where
$K\in\{\R,\C\}$ and $0<\rho\leq1$.  A metric weight is
additive: for a local frame $e$,
\begin{equation}
 \phi_\rho(e):=-\log\norm e_\rho,
 \qquad
 \phi_\rho(ue)=\phi_\rho(e)-\log\abs u_\rho.
\label{eq:metric-transition}
\end{equation}
Its unpowered metric is
\begin{equation}
 \overline\phi(e):=\rho^{-1}\phi_\rho(e),
 \qquad
 \norm e_{\overline\phi,\infty}:=\norm e_{\phi_\rho,\rho}^{1/\rho}.
\label{eq:metric-unpower}
\end{equation}
Division of~\eqref{eq:metric-transition} by $\rho$ gives the usual metric
transition law.

\begin{lemma}
\label{lem:PS-transport}
Under the canonical homeomorphism from a usual complex or real
analytification to the $K_\rho$-analytification, the restricted fibre metric
of \cite[Proposition~3.2]{PS23Hybrid} is transported to the usual complex
metric by \eqref{eq:metric-unpower}.  The fibrewise Monge--Amp\`ere measure
of \cite[Definition~4.8]{PS23Hybrid} is therefore the Boucksom--Eriksson
Archimedean measure of the unpowered metrics
\cite[formula~(8.1) and Proposition~8.3]{BE21}.  This is a consequence of
the stated fibre convention, not a new formal definition of a symbol
$\ddc_\rho$.
\end{lemma}

\begin{proof}
The exponent is explicit in \cite[Proposition~3.2]{PS23Hybrid}.  In the
hybrid-disc notation, put
\[
 \rho_t:=\frac{\log r}{\log\abs t}\quad(0<\abs t\leq r<1).
\]
For the canonical homeomorphism $\beta_t:X_t\to X_t^{\mathrm{hyb}}$, their
proof gives, for every local regular function $f$,
\begin{equation}
 \abs{f(\beta_t(z))}=\abs{f(z)}_\infty^{\rho_t},
 \qquad
 \norm{s(z)}_{\phi_t}
 =\norm{s(\beta_t(z))}_{\phi}^{1/\rho_t}.
\label{eq:PS32-transport}
\end{equation}
Equivalently,
$\log\norm{s}_{\phi_t}/\log_r\abs t$
extends as the hybrid logarithmic norm.  Thus the raw hybrid fibre weight is
$\rho_t$ times the ordinary complex weight, exactly as in
\eqref{eq:metric-unpower}.

For an arbitrary Archimedean $x\in\cM(A)$, the character
\eqref{eq:character-power} gives the same canonical map: on every affine
chart it sends a complex point $z$ to
$f\mapsto\abs{\chi_x(f)(z)}_\infty^{\rho(x)}$.  Formula
\eqref{eq:metric-unpower} is therefore the unique metric transport compatible
with local frames and this seminorm map; a different representing character
only conjugates the construction.

This transport also preserves the psh class used in \cite{PS23Hybrid}.  For
a tropical Fubini--Study metric,
\begin{equation}
 \rho^{-1}\frac1m\max_j
 \{\log\abs{s_j}_\rho+a_j\}
 =\frac1m\max_j
 \{\log\abs{s_j}_\infty+a_j/\rho\},
\label{eq:FS-transport}
\end{equation}
which is usual complex psh by \cite[Theorem~2.11]{PS23Hybrid}.
\cite[Proposition~2.21]{PS23Hybrid} and multiplication by the fixed scalar
$1/\rho$ transport uniform limits as well.

Finally, \cite[Definition~4.8]{PS23Hybrid} defines the family measure at $x$
by restricting the metrics to $X_{\mathcal H(x)}^{\mathrm{an}}$ and applying the
valued-field Monge--Amp\`ere operator of
\cite[Section~4.1]{PS23Hybrid}.  On an Archimedean field,
Boucksom--Eriksson first identify the field with $\R$ or $\C$ only
\emph{up to normalization of its absolute value}
\cite[Section~1.1]{BE21}; their Archimedean current is the classical
Bedford--Taylor current on the usual complex space.  Equations
\eqref{eq:PS32-transport}--\eqref{eq:FS-transport} identify the metric
supplied to that current as $\overline\phi$.

Two further source tests make this reading fail-closed.  First,
\cite[Theorem~4.11]{PS23Hybrid} writes the Archimedean members of a hybrid
family as $(\ddc\phi_t)^n$, where \cite[Proposition~3.2]{PS23Hybrid} has
already defined $\phi_t=\rho_t^{-1}\phi|_{X_t^{\mathrm{hyb}}}$.  Thus the
displayed complex operator acts on the unpowered metric.

Second, \cite[Theorem~8.14]{BE21} is stated over an arbitrary complete valued
field, and its proof uses
\begin{equation}
 \ddc\log\norm{s}_{\phi}\wedge\delta_X
 =\delta_Z-\ddc\phi\wedge\delta_X.
\label{eq:PL-source}
\end{equation}
For $K_\rho$, both
$\log\norm{s}_{\phi_\rho,\rho}$ and $\phi_\rho$ are $\rho$ times their
unpowered counterparts.  Applying the raw Euclidean operator to those
powered weights would turn the left side into
$\rho\delta_Z-\rho\ddc_\infty\overline\phi$, while the right side demanded
by~\eqref{eq:PL-source} would be
$\delta_Z-\rho\ddc_\infty\overline\phi$.  These differ for
$\rho\ne1$ and a nonzero divisor.  The source identity is restored exactly
by the transported operator
\begin{equation}
 \dc_\rho\phi_\rho:=\dc_\infty(\rho^{-1}\phi_\rho),
 \qquad
 \ddc_\rho\phi_\rho:=\ddc_\infty(\rho^{-1}\phi_\rho).
\label{eq:ddc-unpacking}
\end{equation}
Equation~\eqref{eq:ddc-unpacking} is our explicit unpacking of the source
convention, not a formula printed verbatim in either source.
\cite[Proposition~8.3(ii)]{BE21}, which fixes total mass to the
algebraic intersection number, supplies the independent mass check: the raw
Euclidean operator on $\rho\overline\phi$ would multiply an $n$-fold mass
by $\rho^n$.  Hence \cite[Definition~4.8]{PS23Hybrid} has the transported
reading asserted in the lemma.
\end{proof}

\subsubsection{Powered Archimedean normalization}

Let $X/K_\rho$ be projective of dimension $n$; nilpotents, generic
multiplicities, and several top-dimensional components are allowed.

\begin{theorem}
\label{thm:powered-normalization}
For continuous psh-regularizable metrics $\phi_{i,\rho}$, put
$\overline\phi_i=\rho^{-1}\phi_{i,\rho}$.  Under the canonical
identification of analytic spaces,
\begin{equation}
 \MA_\rho(\phi_{1,\rho},\ldots,\phi_{n,\rho})
 =\MA_\infty(\overline\phi_1,\ldots,\overline\phi_n).
\label{eq:PN1-r3}
\end{equation}
All generic lengths in the fundamental cycle occur on both sides.  If $e$
is a nonzero vector in the Deligne line $I_X(L_0,\ldots,L_n)$, then
\begin{equation}
 \lambda_\rho(e;\phi_{0,\rho},\ldots,\phi_{n,\rho})
 =\rho\lambda_\infty
 (e;\overline\phi_0,\ldots,\overline\phi_n).
\label{eq:PN2-r3}
\end{equation}
Equivalently, the recursively defined local height has exactly one factor
$\rho$ in every dimension.
\end{theorem}

\begin{proof}
We prove \eqref{eq:PN1-r3}.  \Cref{lem:PS-transport} identifies the smooth
complex fibre measure with
\[
 \ddc_\infty\overline\phi_1\wedge\cdots\wedge
 \ddc_\infty\overline\phi_n.
\]
For a general scheme, \cite[formula~(8.1)]{BE21} uses the fundamental cycle
\[
 [X]_n=\sum_Zm_Z[Z],
 \qquad
 m_Z=\operatorname{length}(\mathcal O_{X,\eta_Z}).
\]
The resulting measure is the sum of the component measures multiplied by
$m_Z$.
Applying the transport on each reduced top-dimensional component proves the
smooth statement with all generic lengths.

For continuous psh-regularizable metrics, uniformly approximate by smooth
semipositive nets.  Division by the fixed $\rho$ preserves uniform
convergence.  \cite[Theorem~8.4]{BE21}, as used in
\cite[Theorem~4.3]{PS23Hybrid}, gives the unique weakly continuous extension
and hence
\eqref{eq:PN1-r3}.  Over $\R$, both sources define the measure as the
push-forward from complexification.  Complexification commutes with the
transport and introduces no factor two.
\end{proof}

\begin{lemma}
\label{lem:zero-anchor-r3}
Let $X=\operatorname{Spec}A$ be finite and flat over $K_\rho$, and let
$\phi_{0,\rho}$ be a continuous dpsh metric on $L_0$.  Put
$\overline\phi_0=\rho^{-1}\phi_{0,\rho}$.  Under the algebraic
identification of the Deligne line with the norm line,
\begin{equation}
 \log\norm{\langle s\rangle}_{\langle\phi_{0,\rho}\rangle,\rho}
 =\rho\,
 \log\norm{\langle s\rangle}_{\langle\overline\phi_0\rangle,\infty}
\label{eq:zero-metric-r3}
\end{equation}
for every trivializing section $s$ of $L_0$.  Consequently the additive
Deligne weight and every dimension-zero local height are multiplied by
exactly $\rho$, with all residue degrees and nonreduced Artin lengths.
For $K=\R$, the same assertion uses the real-fibre push-forward convention
and has no factor two.
\end{lemma}

\begin{proof}
The algebraic input is \cite[Example~8.15]{BE21}: in dimension zero
$\langle L_0\rangle=N_{A/K}(L_0)$ and
$\langle us\rangle=N_{A/K}(u)\langle s\rangle$.
\cite[Proposition~6.1.1(3)]{Ducrot97CubeIntersection} fixes the positive
algebraic norm orientation, but
neither algebraic statement fixes the multiplicative constant of the
metric.

That constant is fixed by
\cite[Theorem~8.16, formula~(8.11)]{BE21}.  Reconstructed in the present
norm notation, the formula is
\begin{equation}
 \log\norm{\langle s\rangle}_{\langle\phi_0\rangle}
 =\int_{X^{\mathrm{an}}}\log\norm{s}_{\phi_0}\,\delta_X.
\label{eq:BE811-r3}
\end{equation}
Write $A=\prod_jA_j$, with $A_j$ local, residue field $K_j$, and Artin
length $m_j$.  \cite[Example~8.2]{BE21} gives, on the corresponding
analytic points $x_j$,
\begin{equation}
 \delta_X=\sum_j(\dim_KA_j)\delta_{x_j}
 =\sum_jm_j[K_j:K]\,\delta_{x_j}.
\label{eq:zero-cycle-r3}
\end{equation}
Under the canonical identification of the powered and unpowered analytic
spaces, this measure is unchanged, while metric unpowering gives
$\log\norm{s(x_j)}_{\phi_{0,\rho},\rho}
 =\rho\log\norm{s(x_j)}_{\overline\phi_0,\infty}$.
Applying~\eqref{eq:BE811-r3} twice and then~\eqref{eq:zero-cycle-r3}
therefore yields
\begin{align*}
 \log\norm{\langle s\rangle}_{\langle\phi_{0,\rho}\rangle,\rho}
 &=\sum_jm_j[K_j:K]\,
   \log\norm{s(x_j)}_{\phi_{0,\rho},\rho}\\
 &=\rho\sum_jm_j[K_j:K]\,
   \log\norm{s(x_j)}_{\overline\phi_0,\infty}\\
 &=\rho\,
   \log\norm{\langle s\rangle}_{\langle\overline\phi_0\rangle,\infty}.
\end{align*}
This proves~\eqref{eq:zero-metric-r3} and visibly preserves both the residue
degree $[K_j:K]$ and the Artin length $m_j$.  Frame compatibility follows
from \cite[Theorem~8.16, formula~(8.12)]{BE21}, namely
$\int\log\abs u\,\delta_X=\log\abs{N_{A/K}(u)}$; after powering, both
sides acquire the same factor $\rho$.  Finally,
\cite[Remark~8.1]{BE21} defines the real measure as the push-forward from
complexification.  Thus the
degree already present in~\eqref{eq:zero-cycle-r3} is retained and no
additional factor two appears.
\end{proof}

\medskip
\begin{proof}
We prove \eqref{eq:PN2-r3}.  Define the candidate powered Deligne weight by
\begin{equation}
 \widetilde\lambda_\rho
 (e;\phi_{0,\rho},\ldots,\phi_{n,\rho})
 :=\rho\lambda_\infty
 (e;\overline\phi_0,\ldots,\overline\phi_n).
\label{eq:PN2-candidate-r3}
\end{equation}
If $e'=ue$, then
$
\widetilde\lambda_\rho(e')=\widetilde\lambda_\rho(e)
-\rho\log\abs u_\infty
=\widetilde\lambda_\rho(e)-\log\abs u_\rho
$, so this is a metric for the powered absolute value.
At fixed $\rho$, unpowering is additive under tensor products; hence
symmetry and multilinearity of the usual Deligne metric give the same
properties for $\widetilde\lambda_\rho$.

Let $s\in H^0(X,L_0)$ be regular, with divisor $Z$.  The usual additive
restriction formula is
\begin{align*}
 \lambda_\infty(e;\overline\phi_0,\ldots,\overline\phi_n)
 ={}&\lambda_\infty(e|_Z;
       \overline\phi_1|_Z,\ldots,\overline\phi_n|_Z)\\
 &-\int_X\log\norm{s}_{\overline\phi_0,\infty}
       \,\MA_\infty(\overline\phi_1,\ldots,\overline\phi_n).
\end{align*}
Multiplication by $\rho$, the identity
$
 \log\norm{s}_{\phi_{0,\rho},\rho}
 =\rho\log\norm{s}_{\overline\phi_0,\infty}
$, and \eqref{eq:PN1-r3} turn this into the powered restriction formula.  The
dimension-zero lemma supplies the base case, and the uniqueness clause in
\cite[Theorem~8.16]{BE21} identifies the candidate with the target Deligne
metric.

The same calculation is the local-height induction: the restricted height
has one factor $\rho$, while the integral has one factor from its Green
function and none from its measure.  No new factor is accumulated at the
next dimension.
\end{proof}

\subsubsection{A fibre-uniform Chern--Levine--Nirenberg estimate}

For every $0\leq i\leq n$, let
$\phi_{i,\rho},\phi'_{i,\rho}$ be continuous psh-regularizable
semipositive metrics on $L_i$, and put
$f_\rho=\phi_{0,\rho}-\phi'_{0,\rho}$.  In particular, the slot-zero
test is a difference of two metrics in the semipositive domain of the CLN
inequality; it is not an arbitrary continuous function.  Define
\begin{equation}
 C:=
 \begin{cases}
 0,&n=0,\\[2mm]
 2\displaystyle\max_{1\leq i\leq n}
 (L_0L_1\cdots\widehat{L_i}\cdots L_n),&n\geq1.
 \end{cases}
\label{eq:CLN-C-r3}
\end{equation}

\begin{lemma}
\label{lem:CLN-r3}
Assume explicitly that both metrics in every pair
$\phi_{i,\rho},\phi'_{i,\rho}$, including $i=0$, are continuous
psh-regularizable and semipositive.  For $0<\rho\leq1$,
\begin{align}
 &\left|\int_Xf_\rho\,
 \MA_\rho(\phi_{1,\rho},\ldots,\phi_{n,\rho})
 -\int_Xf_\rho\,
 \MA_\rho(\phi'_{1,\rho},\ldots,\phi'_{n,\rho})\right|\notag\\
 &\hspace{30mm}\leq
 C\sum_{i=1}^n
 \norm{\phi_{i,\rho}-\phi'_{i,\rho}}_\infty.
\label{eq:CLN-r3}
\end{align}
The constant is independent of $\rho$.
\end{lemma}

\begin{proof}
If $n=0$, there are no metric slots to change, the two integrals coincide,
and both sides are zero by definition~\eqref{eq:CLN-C-r3}.

Assume $n\geq1$ and first suppose all metrics are smooth semipositive.
Put $\overline\phi_i=\rho^{-1}\phi_{i,\rho}$ and
$\overline\phi'_i=\rho^{-1}\phi'_{i,\rho}$ for every $0\leq i\leq n$,
and $\overline f=\overline\phi_0-\overline\phi'_0$.
The convention transport preserves semipositivity, and
\eqref{eq:PN1-r3} makes the left side equal to
\[
 \rho\left|\int_X\overline f\,
 \MA_\infty(\overline\phi_1,\ldots)
 -\int_X\overline f\,
 \MA_\infty(\overline\phi'_1,\ldots)\right|.
\]
\cite[Lemma~8.6]{BE21} applies because the slot-zero pair
$(\overline\phi_0,\overline\phi'_0)$ is also semipositive.  It gives one
inequality with constant~\eqref{eq:CLN-C-r3}.  Keep that test pair fixed and
interchange primed and unprimed metrics only in slots $1,\ldots,n$ to get
the opposite inequality.  Therefore
\begin{align*}
 \rho\left|\int\overline f\,\MA_\infty(\overline\phi_i)
 -\int\overline f\,\MA_\infty(\overline\phi'_i)\right|
 &\leq\rho C\sum_i
 \norm{\overline\phi_i-\overline\phi'_i}_\infty\\
 &=C\sum_i
 \norm{\phi_{i,\rho}-\phi'_{i,\rho}}_\infty.
\end{align*}
This is the two-sided cancellation: the test function contributes $\rho$,
and unpowering a metric difference contributes $1/\rho$.

For the stated continuous domain, uniformly approximate all
$2(n+1)$ metrics, including the slot-zero pair, by smooth semipositive
metrics.  \cite[Theorem~8.4]{BE21} gives weak convergence of the
Monge--Amp\`ere measures, while uniform convergence of the slot-zero
differences passes the integrals to the limit.  The same constant $C$
therefore survives.  This argument supplies no estimate for an arbitrary
continuous test function outside a difference of two such semipositive
metrics.
\end{proof}

\medskip
\begin{corollary}
\label{cor:p171-production-family-CLN}
Let $X/A$ be flat and projective, let the line bundles above come from $X$,
and assume that for every $0\leq i\leq n$ both metrics in the pair
$\phi_i,\phi'_i$ are continuous psh-regularizable and semipositive on the
fibres.  The same constant works on every fibre for the test
$\phi_0-\phi'_0$.  It is not asserted for an arbitrary continuous test
function.  At $\rho(x)>0$ this is \cref{lem:CLN-r3}; at
$\rho(x)=0$ it is \cite[Lemma~4.10]{PS23Hybrid} and its continuous
semipositive extension.  No division by $\rho$ occurs on the endpoint
locus.
\end{corollary}

\begin{proof}
For $n=0$, $C=0$.  For $n\geq1$, the constant is made only from mixed
algebraic intersection numbers, which are fibre-independent by flatness.
\cref{lem:CLN-r3} has no residual power parameter.  On the
non-Archimedean locus, \cite[Lemma~4.10]{PS23Hybrid} starts with FS metrics in
every slot, including the slot-zero pair.  Uniform FS approximation of all
these semipositive metrics and the same inequality pass the estimate to the
stated continuous domain.
\end{proof}
